% !TITLE 两汉
% !SUMMARY 乐府与古诗的交响
% !ORDER 0

\begin{intro}
两汉四百年，是中国诗歌从集体歌唱走向个人抒情的关键转折。

汉乐府继承了《诗经》"饥者歌其食，劳者歌其事"的现实主义精神，却把镜头对准了更广阔的社会图景。这里有十五从军征的老兵，有饮马长城窟的思妇，有孤儿、有病妇——那些沉默的大多数，第一次在诗歌中发出了自己的声音。乐府诗的语言是质朴的，甚至粗粝的，但正是这种不事雕琢的真诚，让它们穿越两千年依然动人。

而古诗十九首则代表了另一条路径。它们大约产生于东汉末年，作者姓名已经佚失，却标志着文人五言诗的成熟。"行行重行行，与君生别离"——个人的离愁别绪、生命意识、人生无常的感慨，第一次成为中国诗歌的核心主题。这些诗里的主人公是孤独的：他站在高楼上，他走在漫漫长路上，他在庭中折花，他在江畔采芙蓉。这种孤独不是屈原式的家国之忧，而是更普遍、更日常的生命体验。

从乐府的"感于哀乐，缘事而发"到古诗的"婉转附物，怊怅切情"，两汉诗歌完成了从集体到个人、从叙事到抒情的转变。这个转变，为魏晋南北朝的文学自觉铺好了道路。
\end{intro}
